\section{数列}
\label{sec:analysis-sequence-basics}

\subsection{数列的定义}

\begin{definition}
给每个正整数\(n\)指定唯一的实数\(a_n\)，按下标顺序排列得到
\begin{equation*}
 a_1,a_2,\ldots,a_n,\ldots,
\end{equation*}
称为一个{\heiti{}数列}，记作\(\{a_n\}\)。\(a_n\)是数列的第\(n\)项，\(n\)称为下标。若只取下标\(1,2,\ldots,N\)，便得到有限数列；本章讨论极限时，均指无穷数列。
\end{definition}

数列中的数可以重复，排列顺序也有意义。例如，\(1,-1,1,-1,\ldots\)与\(-1,1,-1,1,\ldots\)具有相同的取值集合，却是不同的数列。

若用关于\(n\)的式子直接给出\(a_n\)，称为数列的{\heiti{}通项公式}，例如\(a_n=2n-1\)。有了通项公式，就可以由下标确定相应的项。

记前\(n\)项的和为\(S_n\)，即
\begin{equation}\label{eq:analysis-partial-sum-definition}
 S_n=\sum_{k=1}^{n}a_k=a_1+a_2+\cdots+a_n.
\end{equation}
其中\(k\)只是求和下标，可以换成其他字母。约定\(S_0=0\)，则
\begin{equation}\label{eq:analysis-term-from-sum}
 a_n=S_n-S_{n-1},\qquad n\ge1.
\end{equation}
前\(n\)项的积类似地记作\(\prod_{k=1}^{n}a_k\)；特别地，\(n!=1\cdot2\cdots n\)，并约定\(0!=1\)。

\begin{example}
写出数列\(a_n=(-1)^{n-1}n\)的前五项。若另一数列的前\(n\)项和为\(S_n=n(n+1)\)，求它的通项。
\end{example}

只知道有限多项，通常不能唯一确定后面的项。例如，\(a_n=n\)与\(b_n=n+(n-1)(n-2)(n-3)\)的前三项相同，第四项却不同。观察可以帮助猜测通项，证明还须使用题目给出的完整条件。

\begin{example}
分别写出\(1,0,1,0,\ldots\)与\(1,-1/2,1/3,-1/4,\ldots\)的一个通项公式。对数列\(1,1/2,1/2,1/3,1/3,1/3,\ldots\)，其中\(1/k\)依次重复\(k\)次，试用下标的取值范围表示它的通项。
\end{example}

\subsection{单调性与有界性}

若对所有\(n\)都有\(a_{n+1}\ge a_n\)，称\(\{a_n\}\){\heiti{}单调递增}；将不等号换为\(\le\)，便是{\heiti{}单调递减}。取严格不等号时，称为严格单调。证明单调性通常考察差\(a_{n+1}-a_n\)的符号；对正数列，也可比较比值\(a_{n+1}/a_n\)与\(1\)的大小。

若有常数\(M\)，使所有项都满足\(a_n\le M\)，称数列有上界；下界类似。若上下界都有，称数列{\heiti{}有界}，也就是存在\(K>0\)，使\(|a_n|\le K\)对所有\(n\)成立。

有上界不表示上界一定是某一项。例如\(a_n=n/(n+1)\)严格递增，以\(1\)为上界，却没有最大项。数列是否单调、是否有界，将是后面判断收敛性的重要依据。

\begin{example}
判断\(a_n=(2n+1)/(n+1)\)、\(b_n=(-1)^n(1-1/n)\)、\(c_n=n(6-n)\)的单调性与有界性，并求它们存在的最大项、最小项。有界与单调能否互相推出？
\end{example}

\subsection{数学归纳法}

有些命题涉及每个正整数，逐个代入不能穷尽。若能从第一步出发，证明成立的一步总能推出下一步，就可以覆盖所有正整数。

\begin{theorem}[第一数学归纳法]
设\(P(n)\)是关于正整数\(n\)的命题。若\(P(1)\)成立，且对任意\(k\ge1\)，都有\(P(k)\Rightarrow P(k+1)\)，则\(P(n)\)对所有正整数成立。
\end{theorem}

归纳步骤暂时假设\(P(k)\)成立，是为了证明一个蕴含关系；它与已经验证的起始情形共同起作用。若命题从\(n_0\)开始，只须相应改变起点。

例如，\(1+3+\cdots+(2n-1)=n^2\)在\(n=1\)时成立。假设对\(n=k\)成立，再添下一项\(2k+1\)，便得\(k^2+2k+1=(k+1)^2\)，故公式对所有正整数成立。

\begin{example}
用数学归纳法证明\(\sum_{k=1}^n k(k+1)=n(n+1)(n+2)/3\)。再对\(x\ge-1\)证明Bernoulli不等式\((1+x)^n\ge1+nx\)，说明归纳步骤中为何可以将不等式两边乘以\(1+x\)。
\end{example}

有时，证明下一步需要使用前面若干步的结论，所用的下标也未必紧邻下一步。这时可以采用下面的形式。

\begin{theorem}[第二数学归纳法]
\label{thm:analysis-strong-induction}
设\(n_0\)为正整数，\(P(n)\)是关于整数\(n\ge n_0\)的命题。若\(P(n_0)\)成立，且对任意\(k\ge n_0\)，由\(P(n_0),P(n_0+1),\ldots,P(k)\)同时成立能够推出\(P(k+1)\)，则\(P(n)\)对所有\(n\ge n_0\)成立。
\end{theorem}

\begin{proof}
将“\(P(n_0),P(n_0+1),\ldots,P(k)\)均成立”记为\(Q(k)\)。由条件，\(Q(n_0)\)成立；若\(Q(k)\)成立，则\(P(k+1)\)成立，连同已有的结论便得到\(Q(k+1)\)。对\(Q(k)\)使用第一数学归纳法，即得结论。
\end{proof}

第二数学归纳法也称为{\heiti{}强归纳法}，归纳假设允许使用从起点到第\(k\)步的全部结论。两种形式可以互相推出，使用时依照证明中实际需要的结论选择即可。

\begin{example}
用第二数学归纳法证明：每个大于\(1\)的整数都可以写成若干个质数的乘积。
\end{example}

对于用到前两项的递推关系，可以先验证前两个起始情形，再同时使用这两项的结论作归纳推导；也可以将相邻两个命题合成一个命题，使用第一数学归纳法。起始情形需要验证几项，应当由递推所需的条件决定。

\begin{exercise}
\label{ex:analysis-induction}
用适当的数学归纳法完成下列各题。

（1）证明\(1^3+2^3+\cdots+n^3=[n(n+1)/2]^2\)。

（2）证明\(10^n-9n-1\)能被\(81\)整除。

（3）证明当\(n\ge4\)时，\(2^n\ge n^2\)。

（4）设\(a_1=1,a_2=2,a_{n+2}=a_{n+1}+a_n\)。证明\(n\le a_n\le2^{n-1}\)。

（5）设\(a_1=1\)，且当\(n\ge2\)时，\(a_n=\sum_{k=1}^{n-1}a_k\)。猜测并证明通项公式，再将给定关系化为相邻两项之间的递推关系。

（6）证明每个正整数都能表示为若干个互不相同的\(2\)的幂之和，其中指数为非负整数。
\end{exercise}

